% fig_plddt.tex -- per-residue confidence trace of the aRES monomer model.
% Data: plddt.dat, generated from the B-factor column of aRES_Zn_segA.pdb
% (per-residue mean). That column is a pLDDT-like confidence field, verified
% against the per-region means recorded in docs/07: HIGH = confident.
\begin{tikzpicture}
\begin{axis}[
    width=0.95\textwidth, height=5.2cm,
    xlabel={residue}, ylabel={pLDDT},
    xmin=1, xmax=447, ymin=20, ymax=100,
    xtick={1,50,100,150,200,250,300,350,400,447},
    ytick={20,40,60,80,100},
    tick label style={font=\footnotesize},
    label style={font=\small},
    axis line style={black!70},
    grid=major, grid style={black!12, line width=0.3pt},
    clip mode=individual,
]

  % confidence bands
  \addplot[draw=none, fill=green!12, forget plot]
    coordinates {(1,70) (447,70) (447,100) (1,100)} \closedcycle;
  \addplot[draw=none, fill=red!10, forget plot]
    coordinates {(1,20) (447,20) (447,50) (1,50)} \closedcycle;

  % the two low-confidence regions called out in docs/07
  \addplot[draw=none, fill=red!28, forget plot]
    coordinates {(121,20) (125,20) (125,100) (121,100)} \closedcycle;
  \addplot[draw=none, fill=orange!22, forget plot]
    coordinates {(392,20) (408,20) (408,100) (392,100)} \closedcycle;

  \addplot[black!80, line width=0.6pt] table[x=resi, y=plddt] {figures/plddt.dat};

  \node[font=\scriptsize, anchor=south, text=red!65!black, align=center]
    at (axis cs:123,100) {hinge};
  \node[font=\scriptsize, anchor=south, text=orange!65!black, align=center]
    at (axis cs:400,100) {loop};

\end{axis}
\end{tikzpicture}
